% ============================================
% Chapter 1: State of the Art
% ============================================

\chapter{State of the Art}\label{ch:state_of_art}

\section{Introduction}

Network intrusion detection has evolved from manually maintained signature databases to data-driven machine learning systems capable of recognising behavioural anomalies. This chapter reviews the foundations of port scan detection, compares the principal families of detection algorithms, and introduces the active defence concepts that underpin the AEGIS Entropy architecture. The objective is to position the project within the current research and industrial landscape, and to justify the technical choices made in the subsequent chapters.

The chapter is organised as follows. Section~\ref{sec:portscan} defines port scanning and presents a taxonomy of common techniques. Section~\ref{sec:ids} contrasts classical and AI-based intrusion detection approaches. Section~\ref{sec:ml_dl} compares Machine Learning and Deep Learning for intrusion detection and explains why Machine Learning was selected for this project. Section~\ref{sec:supervised_unsupervised} discusses supervised and unsupervised learning paradigms. Sections~\ref{sec:mtd} and~\ref{sec:honeypot} introduce Moving Target Defence and honeypot-based deception. Section~\ref{sec:soa_conclusion} concludes with the positioning of AEGIS Entropy.

\section{Port Scan Attack Detection}\label{sec:portscan}

\subsection{Definition}

Port scanning is the systematic probing of a host or network to identify open ports, active services, and potential vulnerabilities. It is almost always the first reconnaissance step of a cyberattack: before an adversary can exploit a service, they must know that the service exists and understand its behaviour. Port scans generate distinctive traffic patterns---half-open connections, high volumes of rejected packets, and unusual inter-arrival times---that can be distinguished from normal host-to-host communication when the right features are monitored.

From a machine learning perspective, port scan detection is a binary classification problem: given a set of flow-level or packet-level features, assign the label \textsc{Benign} or \textsc{PortScan}. The difficulty lies in the diversity of scanning strategies, the similarity of some scan patterns to legitimate behaviour, and the need to operate in real time.

\subsection{Taxonomy of Port Scan Techniques}

Table~\ref{tab:scan_taxonomy} summarises the most common port scan techniques. Each technique produces a different signature in terms of TCP flags, completion state, and response behaviour, which in turn determines the features most useful for detection.

\begin{table}[H]
\centering
\caption{Classification of port scan techniques}
\label{tab:scan_taxonomy}
\begin{tabularx}{\textwidth}{lX}
\toprule
\textbf{Scan Type} & \textbf{Description} \\
\midrule
SYN Scan & Half-open scan using SYN packets without completing the TCP handshake. Fast and stealthy because it avoids full connection logs. \\
TCP Connect & Full TCP connection establishment to each target port. Easier to detect but works without raw socket privileges. \\
UDP Scan & Sends UDP packets to detect open UDP services. Relies on ICMP "port unreachable" responses for closed ports. \\
XMAS Scan & Sets FIN, PSH, and URG flags to elicit responses from closed ports. Often blocked by modern stacks but still used for firewall mapping. \\
ACK Scan & Sends ACK packets to map firewall rules by analysing RST responses. Useful for determining whether a port is filtered. \\
Service Detection & Probes identified open ports for service version information after the initial scan. \\
OS Detection & Analyses TCP/IP stack behaviour to fingerprint the operating system of the target. \\
\bottomrule
\end{tabularx}
\end{table}

SYN scans and UDP scans are the most common in real-world reconnaissance, while XMAS and ACK scans are often used to evade simple stateless firewalls. Service and OS detection occur in later phases and produce longer, more structured sessions.

\section{Network Intrusion Detection Systems}\label{sec:ids}

\subsection{Signature-Based Detection}

Signature-based Intrusion Detection Systems (IDS) such as Snort and Suricata compare observed traffic against databases of known attack patterns. They are highly accurate when the attack signature is known and well-defined, and they provide clear, explainable alerts. However, they require continuous manual updates, cannot detect zero-day or polymorphic attacks, and often generate many alerts in environments with diverse legitimate traffic.

\subsection{Anomaly-Based Detection}

Anomaly-based IDS build a statistical or behavioural model of normal network activity and flag significant deviations. This approach can detect previously unseen attacks, including slow scans and custom reconnaissance tools. The main challenge is defining ``normal'' behaviour: legitimate administrative tools, vulnerability scanners, and cloud orchestration can all appear anomalous, leading to false positives.

\subsection{Classical IDS/IPS vs AI-Based IDS/IPS}

Artificial Intelligence, and more specifically machine learning, offers a fundamentally different approach. Rather than relying on manually curated rules, AI-based systems learn statistical models of normal and malicious behaviour from labelled traffic datasets such as CICIDS2017~\cite{cicids2017}. Table~\ref{tab:ids-comparison} contrasts the two paradigms.

\begin{table}[H]
\centering
\caption{Classical IDS/IPS vs.\ AI-Based IDS/IPS}
\label{tab:ids-comparison}
\begin{tabularx}{\textwidth}{XX}
\toprule
\textbf{Classical IDS/IPS} & \textbf{AI-Based IDS/IPS} \\
\midrule
Relies on predefined signatures & Learns behavioural patterns from data \\
Fails against unknown/novel attacks & Detects anomalies beyond known signatures \\
Requires constant manual updates & Self-adapts through model retraining \\
No autonomous response capability & Supports automated threat neutralisation \\
High false positive rate at scale & Optimised precision/recall via ML tuning \\
Static and network-topology-specific & Generalisable across diverse network types \\
\bottomrule
\end{tabularx}
\end{table}

AI-based detection is not a replacement for classical defences; it is a complementary layer that reduces alert fatigue and improves coverage against evolving reconnaissance techniques.

\section{Machine Learning vs Deep Learning for Intrusion Detection}\label{sec:ml_dl}

\subsection{Deep Learning Approaches}

Deep Learning approaches for intrusion detection include Convolutional Neural Networks (CNNs), Recurrent Neural Networks (RNNs) and Long Short-Term Memory networks (LSTMs), and autoencoders. CNNs have been applied to raw packet bytes or to flow images generated from traffic features; RNNs and LSTMs model temporal dependencies in packet sequences; autoencoders learn compressed representations of normal traffic and flag reconstructions with high error as anomalies.

These methods can achieve high accuracy on large datasets and can discover features automatically. However, they require substantial training data, significant computational resources, and long hyperparameter tuning cycles. They also tend to be ``black boxes'' that are difficult to explain to SOC analysts, which is a serious drawback for security-critical deployments.

\subsection{Machine Learning Approaches}

Machine Learning approaches for intrusion detection include tree-based ensembles (Random Forest~\cite{breiman2001}, XGBoost~\cite{chen2016}), Support Vector Machines, k-Nearest Neighbours, and Na\"{i}ve Bayes classifiers. These methods operate directly on structured, hand-crafted features such as packet counts, flag statistics, and inter-arrival times. They train quickly, provide feature importance scores, and are robust on moderate-sized datasets such as the CICIDS2017 PortScan subset used in this project.

\subsection{The Choice of Machine Learning for AEGIS Entropy}

AEGIS Entropy uses Machine Learning rather than Deep Learning for three reasons:

\begin{enumerate}
    \item \textbf{Data nature:} The CICIDS2017 dataset provides structured, flow-level features. Tree-based models are state-of-the-art for tabular data of this kind.
    \item \textbf{Interpretability:} SOC analysts need to understand why an alert was raised. Random Forest and XGBoost provide feature importance scores that explain which flow characteristics triggered the detection.
    \item \textbf{Efficiency:} The project must run in real time on a virtualised lab environment. Tree-based models offer low-latency inference with modest hardware requirements.
\end{enumerate}

\section{Supervised vs Unsupervised Learning}\label{sec:supervised_unsupervised}

\subsection{Supervised Learning}

Supervised learning uses labelled training data to learn a mapping from input features to class labels. In AEGIS Entropy, Random Forest and XGBoost are trained on CICIDS2017 flow records labelled \textsc{Benign} or \textsc{PortScan}. The models learn decision boundaries that separate the two classes and can then classify new, unseen flows. Supervised models excel when the attack patterns in production traffic resemble those seen during training.

\subsection{Unsupervised Learning}

Unsupervised learning does not require labelled data. Instead, it learns a model of normal behaviour and flags deviations as anomalies. The Isolation Forest algorithm~\cite{liu2008} is used in AEGIS Entropy as a complementary layer for detecting slow or unknown scan patterns. Isolation Forest isolates anomalies by randomly selecting features and split values; anomalies are ``few and different'' and therefore require fewer splits to isolate.

In practice, Isolation Forest must be trained primarily on benign traffic to be effective. When trained on a dataset where attacks are the majority---as in the CICIDS2017 PortScan slice---its assumptions are violated, a phenomenon analysed in detail in Chapter~\ref{ch:model_training}.

\section{Moving Target Defense}\label{sec:mtd}

Moving Target Defence (MTD) is a proactive security strategy that continuously changes the attack surface of a system to invalidate reconnaissance data collected by adversaries. Rather than passively detecting and blocking attacks, MTD introduces uncertainty: a port mapping, IP address, or service fingerprint obtained by an attacker becomes stale within seconds or minutes.

Common MTD techniques include:

\begin{itemize}
    \item \textbf{Port mutation:} Dynamically rotating the listening ports of protected services.
    \item \textbf{IP hopping:} Periodically changing host or service IP addresses.
    \item \textbf{Fingerprint randomisation:} Altering TCP/IP stack signatures to mislead OS and service fingerprinting tools.
\end{itemize}

AEGIS Entropy adopts port mutation and deception-based redirection rather than full IP hopping, because port-level changes are sufficient to disrupt reconnaissance while remaining practical in a controlled lab environment. MTD is discussed in the context of NIST guidance~\cite{nist2020} and the broader MTD literature~\cite{jajodia2011}.

\section{Honeypot and Deception-Based Defense}\label{sec:honeypot}

A honeypot is a deliberately exposed decoy system that mimics a legitimate service while containing no real operational data. Any interaction with a honeypot is, by definition, unauthorised and constitutes a high-confidence threat signal. Honeypots can be classified by interaction level:

\begin{itemize}
    \item \textbf{Low-interaction honeypots} emulate services at the application layer. They are lightweight and easy to deploy but offer limited attacker engagement.
    \item \textbf{Medium-interaction honeypots} provide more realistic sessions, allowing limited command execution and deeper behavioural observation.
    \item \textbf{High-interaction honeypots} are full operating systems or services. They capture the richest attacker behaviour but require careful isolation and higher maintenance.
\end{itemize}

In AEGIS Entropy, lightweight decoy listeners are deployed alongside the real services. When a scanner contacts a decoy port, the event is logged and can be used as an additional feature for the ML model, increasing detection confidence.

\section{Conclusion}\label{sec:soa_conclusion}

The state of the art shows a clear trajectory from static signatures to adaptive machine learning detection, and from passive alerting to active defence. AEGIS Entropy combines these trends into a single architecture: supervised learning (Random Forest and XGBoost) for known scan patterns, unsupervised learning (Isolation Forest) for behavioural anomalies, and active defence (MTD and honeypot redirection) to increase attacker cost and collect threat intelligence. Chapter~\ref{ch:understanding} translates these foundations into the specific problem statement, user requirements, and data dictionary for the AEGIS Entropy system.
